% ============================================================
% INTERPRETAÇÃO SUGERIDA — Figura 9 (REVISADA v2)
% Pearson |r| entre atributos dos campos dinâmicos e termos do LEC (EPALL)
%
% Este arquivo contém o texto sugerido para a subseção
% "Relationship between energy terms and dynamical fields"
% (Seção \ref{sec:correlation} do artigo).
%
% REVISÃO v2 — 2026-04-30 — Correções:
%   1. Convencão de sinal do AdvT_850: setor norte e oeste têm advecção FRIA
%      (NW = -60e-6 K/s no compósito); setor leste tem advecção quente (+33e-6).
%      A interpretação original ("advecção quente a norte e oeste") estava incorreta.
%   2. Convenção de sinal do Ck: Ck < 0 = instabilidade barotropica (fluxo zonal →
%      perturbação), conforme equações do artigo (∂Ke/∂t = +Ck) e legenda da Fig. 5
%      ("negative values = mean flow to eddy"). Todos os centroides dos clusters têm
%      Ck < 0. A interpretação anterior estava com o sinal invertido.
%   3. PV_200: todos os valores são NEGATIVOS no HS. Contraste E-W > 0 SIGNIFICA
%      que o LESTE é menos negativo (anticiclônico) e o OESTE é mais negativo
%      (ciclônico = cavado). Portanto, a configuração PV_200 E-W = leste anticiclônico
%      + oeste ciclônico. A interpretação anterior ('cavado a leste') estava errada.
%      - r(Ca/Ae, border_west PV_200) < 0: mais ciclônico a oeste → MAIS Ca/Ae
%        (fortalece, não suprime, a conversão barotrópica).
%   4. Nota adicionada: r_abs(AdvT_850 × KE_adv_250) = +0.607 (compartilhamento de
%      forçantes advectivas — ponto levantado pelo Danilo).
%   5. PV_850 excluído da análise detalhada (7,7% dos pares acima do limiar).
%   6. Parágrafo "primera vez" removido (era desnecessário).
%   7. Nota sobre adição de Spearman ρ nos métodos incluída.
%
% Convenção de sinal do Ck confirmada via:
%   - Equações do artigo: ∂K_Z/∂t = ... − Ck + ...; ∂K_E/∂t = ... + Ck + ...
%     → Ck > 0 implica K_E ganha energia (Kz→Ke = INSTABILIDADE barotropica)
%   - No entanto, a legenda da Fig. 5 diz "negative values = mean flow to eddy"
%     para os SUBtermos. Os centroides dos clusters confirmam que EP1 (o mais
%     barotropico) tem Ck_int = -12.89 (mais negativo), logo a convenção dos
%     SUBtermos (e do Ck do artigo) é Ck < 0 = instabilidade barotropica.
%   IMPORTATE: Há uma inconsistência entre o texto "C_K signifies the transformation
%   from K_E to K_Z" e as equações diferenciais. Pelo texto e pela legenda da Fig. 5,
%   a convenção utilizada é Ck < 0 = instabilidade barotropica. Confirmar com co-autores.
%
% Método de análise dos valores numéricos:
%   - Valores exatos de r (com sinal) obtidos de:
%       results/lec_field_dependence/step7_predep_absolute_epall.csv
%   - Valores de campo (setor, compósito) obtidos de:
%       data/era5_ep_structure/precomputed_composites_epall.nc
%   - Correlação AdvT × KE_adv obtida de:
%       results/lec_field_dependence/step6_integrated_absolute.csv
%   - N amostral: 2729–2733 ciclones (EPALL, fase de intensificação)
%   - Limiar crítico para significância (p < 0.05, bicaudal): |r| = 0.038
%   - TODAS as correlações exibidas (|r| ≥ 0.20) são altamente significativas.
%
% Gerado em: 2026-04-30
% ============================================================

%% INÍCIO DO TEXTO SUGERIDO %%

\subsection{Relationship between energy terms and dynamical fields}
\label{sec:correlation}

The question of whether the LEC terms are connected to the dynamical fields presented
in Sect.~\ref{sec:composites} naturally arises. Here, we evaluate whether linear
relationships exist between the LEC terms and scalar quantities derived from the
dynamical fields for all cyclones considered in this study, using Pearson correlation
coefficients ($r$). The monotonicity of each association was verified using
Spearman's rank correlation coefficient ($\rho$); sign agreement between Pearson $r$
and Spearman $\rho$ was found for 100\% of the 126 pairs with $|r| \geq 0.20$, and
only two pairs showed a modest magnitude divergence ($|\,|r| - |\rho|\,| > 0.10$).
Although correlation does not necessarily imply causality, it can provide supporting
evidence when the association is physically plausible. This is shown in
Fig.~\ref{fig:9_pearson_epall_by_field_type}, where only values with $|r| > 0.2$ are
displayed. Although this threshold is usually considered to indicate a weak relationship,
the large variability in the cyclone dataset and the complexity of LEC terms make these
associations noteworthy. Given the sample size ($n \approx 2733$), the critical value
for statistical significance at $p < 0.05$ (two-tailed Pearson test) is
$|r_{\mathrm{crit}}| \approx 0.037$; consequently, all displayed associations
($|r| \geq 0.20$) are highly significant ($p < 10^{-21}$ in the worst case). It must
be noted, however, that statistical significance does not preclude the possibility of
spurious associations in a large sample: physical plausibility and sign consistency
with Spearman's $\rho$ are therefore used as primary interpretive filters.

The $\mathrm{PV}_{850}$ field yielded only 7 of 91 possible pairs (7.7\%) above the
display threshold, markedly fewer than any other field analyzed, and its maximum
correlation was $|r| = 0.27$. Given this sparsity and the inconsistency of the
associations with physical mechanistic inference (see below), the
$\mathrm{PV}_{850}$ panel is not discussed in detail here.

Indeed, some correlations reach relatively high values ($|r| > 0.5$), despite the fact
that LEC terms often involve multiple contributions, including departures from the zonal
mean, deviations from the domain mean, and different forms of spatial averaging.

% --- NOTE FOR METHODS SECTION ---
% The following sentence should be added to the statistical analysis paragraph in
% Sect.~\ref{sec:methodology}: "Associations between each canonical LEC term and each
% dynamical-field scalar feature were quantified using both Pearson's correlation
% coefficient $r$ and Spearman's rank correlation $\rho$, the latter providing
% robustness to outliers and non-linear monotone relationships. Only pairs with
% $|r| \geq 0.20$ are shown in Fig.~\ref{fig:9_pearson_epall_by_field_type};
% Spearman $\rho$ values are reported in Table~S? of the supplementary material."
% ---------------------------------

% ---------------------------------------------------------------
% Panel (a): AdvT 850
% ---------------------------------------------------------------

For the $\mathrm{AdvT}_{850}$ field (Fig.~\ref{fig:9_pearson_epall_by_field_type}a),
the highest correlations are found for LEC terms related to APE, as expected from
quasigeostrophic theory. A key aspect of the interpretation is the spatial structure
of temperature advection in Southern Hemisphere extratropical cyclones: in the
storm-centered composite (EPALL), the western and northwestern sectors exhibit
strongly negative temperature advection (cold advection, $\sim -60 \times 10^{-6}$
K\,s$^{-1}$ at NW), consistent with the cold front located to the west and southwest
of the cyclone center. The eastern and northeastern sectors show positive values
(warm advection, $\sim +33 \times 10^{-6}$ K\,s$^{-1}$), corresponding to the warm
sector and warm conveyor belt on the eastern flank \citep{dacre2012extratropical}. The
northern sector has, on average, slightly negative advection, as it samples the cold
air behind the advancing cold front at a slightly equatorward position.

For $C_a$ (baroclinic conversion), the strongest correlations are concentrated in the
western sector ($r = -0.46$, $\rho = -0.55$) and the northern sector ($r = -0.44$,
$\rho = -0.55$), alongside the domain mean ($r = -0.44$). The negative sign is
physically consistent with the above structure: stronger cold advection in the western
and northern sectors (more negative $\mathrm{AdvT}$) indicates a more vigorous cold
front and a steeper meridional temperature gradient, which supports stronger baroclinic
conversion \citep{davies2015quasigeostrophic,de2025lorenz}. The positive correlation
with the absolute domain mean ($r = +0.25$ for $C_a$) confirms the complementary role
of the warm sector: larger-amplitude advection in both sectors (strong frontal contrast)
correlates with stronger $C_a$.

For $A_e$ (eddy available potential energy), correlations are the strongest among all
LEC terms in this panel, reaching $r = -0.55$ ($\rho = -0.63$) for the northern sector
and $r = -0.54$ for the northern border. The negative sign reflects the same mechanism:
stronger cold advection in the cold sector implies a deeper baroclinic zone, which
sustains a larger $A_e$ reservoir. The absolute domain mean correlates positively with
$A_e$ ($r = +0.38$), again reflecting the amplitude of the full frontal contrast. A
notable divergence between Pearson and Spearman for $C_a \times \mathrm{sector\_N}$
($|r| = 0.44$ vs $|\rho| = 0.55$) suggests that the $C_a$–temperature advection
relationship has a moderately non-linear component, consistent with the steep frontal
zone sometimes associated with explosive cyclogenesis.

Significant correlations are also found for $BA_e$ (boundary flux of $A_e$): eastern
border ($r = +0.28$), eastern sector ($r = +0.28$), and E–W contrast ($r = +0.33$),
while the western border and sector show negative correlations ($r = -0.24$ and
$-0.37$, respectively). The physical interpretation is straightforward: warm advection
on the eastern side increases $BA_e$ (Ae import through the eastern boundary), whereas
cold advection on the western side decreases $BA_e$ (Ae export through the western
boundary). Consequently, net Ae import ($BA_e > 0$) requires precisely the
warm-east/cold-west asymmetry that defines the frontal structure of SH cyclones; the
E–W contrast correlation captures both effects simultaneously. Finally, relatively weak
but significant correlations involving $K_e$ ($|r| = 0.22$–$0.35$) are likely
indirect: $C_a$ and $BA_e$ increase the $A_e$ reservoir, which is subsequently
converted into $K_e$ through $C_e$. The marginal correlation for $C_k$ with the domain
mean ($r = +0.20$) is noted in the next paragraph.

% ---------------------------------------------------------------
% Panel (b): AFC 250
% ---------------------------------------------------------------

For the $\mathrm{AFC}_{250}$ field (Fig.~\ref{fig:9_pearson_epall_by_field_type}b),
the highest correlations are found for the reservoir terms $K_e$ and $A_e$, both
associated with the amplitude of AFC ($|\mathrm{AFC}|$ domain mean;
$r_{K_e} = +0.53$, $r_{A_e} = +0.51$). The positive sign indicates that stronger
upper-level ageostrophic geopotential flux convergence (in absolute value, hence
capturing both positive convergence and divergence events) is associated with larger
eddy energy reservoirs. This is physically consistent with the downstream development
framework of \citet{orlanski1991life,orlanski1993ageostrophic}: positive AFC at the
upper-tropospheric level reflects local accumulation of eddy kinetic energy
by ageostrophic energy import, while negative AFC indicates downstream dispersal.
The large absolute-value correlation, rather than the signed mean, implies that the key
signal is the \emph{intensity} of the upper-level energy redistribution, regardless of
its direction, rather than a systematic bias toward convergence or divergence.

Regarding the signed domain-mean AFC, $C_a$ shows $r = +0.43$ and $C_k$ shows
$r = -0.22$, both statistically robust. The positive $C_a$–AFC relation suggests that,
on average, cyclones with more organized upper-level downstream energy redistribution
(positive net AFC, consistent with ageostrophic geopotential convergence near the
cyclone core) \citep{pinheiro2022contributions,pinheiro2024deepening} tend to have
stronger baroclinic conversion. For $C_k$ (with the convention used here:
$C_k < 0$ represents conversion from mean-flow kinetic energy to eddy kinetic energy,
i.e., barotropic instability; see Fig.~5), the negative $C_k$–AFC association
($r = -0.22$ for domain mean, $r = -0.22$ for eastern sector) indicates that cyclones
with stronger barotropic instability (more negative $C_k$) tend to occur in
environments with higher AFC convergence, particularly on the eastern flank. The spatial
breakdown (western border: $r = +0.18$, below threshold) suggests that barotropically
dominant cyclones are associated with a broader, more symmetric AFC pattern, rather
than the predominantly eastward-dispersing AFC signature typical of classically
baroclinic downstream development \citep{orlanski1991life}. This is consistent with the
composite analysis for EP1 (Sect.~\ref{sec:composites}).

The spatial pattern of AFC correlations with $BA_e$ is also noteworthy:
$BA_e$ correlates negatively with the eastern border and the E–W contrast
($r = -0.29$ and $r = -0.41$, respectively), and positively with the western border
($r = +0.38$) and the western sector ($r = +0.20$). This indicates that a stronger
AFC on the western than on the eastern flank — consistent with upper-level energy
convergence upstream of the cyclone center — is associated with more import of eddy
APE into the domain. This is physically interpretable as a consequence of the
large-scale trough–ridge couplet west of the surface cyclone providing an environment
for APE import, as noted in the composite analysis for EP2
(Sect.~\ref{sec:composites}).

% ---------------------------------------------------------------
% Panel (c): KE adv 250
% ---------------------------------------------------------------

For the $\mathrm{KE\_adv}_{250}$ field (Fig.~\ref{fig:9_pearson_epall_by_field_type}c),
the highest individual correlation is $A_e \times \mathrm{sector\_W}$
($r = -0.45$, $\rho = -0.45$), indicating that stronger westward export of kinetic
energy (negative KE advection on the western flank of the domain) is associated with
larger eddy APE. This result may reflect the kinematic relationship between the upper-level
jet structure and the low-level thermal gradient: when the jet produces net kinetic
energy export to the west of the cyclone, the baroclinic zone supporting $A_e$ tends to
be more robust \citep{de2025lorenz}. The amplitude of KE advection ($|\mathrm{KE\_adv}|$
domain mean) is also significantly correlated with $A_e$ ($r = +0.43$) and $C_a$
($r = +0.34$), reinforcing the interpretation that the overall upper-level energy
redistribution efficiency correlates with the baroclinic energy pathway.

Notably, $C_k$ shows positive correlations with several KE advection features:
domain mean ($r = +0.33$), northern sector ($r = +0.26$), and southern and eastern
sectors ($r = +0.24$, $+0.27$). Given that $C_k < 0$ represents barotropic instability,
the positive correlations indicate that 
emph{less} negative $C_k$ (weaker barotropic instability) is associated with
\emph{higher} (more positive) upper-level KE advection into the domain. Equivalently,
cyclones with the strongest barotropic instability (EP1-like, most negative $C_k$)
tend to exhibit more net KE export at 250 hPa. This is consistent with the EP1 mature
and decay stages, which are characterized by large eddy-energy export
(Sect.~\ref{sec:composites}), reflecting a dispersive upper-level wave pattern
during the wave-breaking phase.

An important caveat for the interpretation of KE advection correlations is that
the amplitude of 850-hPa temperature advection and of 250-hPa kinetic energy advection
are themselves strongly coupled: the cross-correlation between their
domain-absolute-mean features is $r = +0.61$ ($\rho = +0.64$, $n = 2733$). This
suggests that a substantial fraction of the $\mathrm{KE\_adv}_{250}$–LEC associations
may reflect shared advective forcing rather than a physically independent upper-level
process. For example, more intense cyclones have both a stronger cold front (large
$|\mathrm{AdvT}_{850}|$) and stronger upper-level KE redistribution (large
$|\mathrm{KE\_adv}_{250}|$), so the correlations of KE advection features with $A_e$
and $C_a$ partially overlap with those already described for $\mathrm{AdvT}_{850}$.

For $G_e$ (diabatic generation of eddy APE), significant negative correlations with KE
advection are found primarily in the southern border ($r = -0.38$, $\rho = -0.43$)
and the southern sector ($r = -0.30$). The negative sign indicates that enhanced
kinetic energy export to the south at 250 hPa — likely reflecting a southward
ageostrophic wind component consistent with an anticyclonic outflow pattern south of the
upper-level trough — is associated with stronger diabatic generation. This is
physically plausible: latent heat release in convective regions embedded in the cyclone
warm sector tends to be organized south of the upper-level jet core, consistent with
upper-tropospheric divergence patterns \citep{huo1999interaction,attinger2021systematic}.

For $BK_e$ (boundary flux of $K_e$), the only correlation above the threshold is with
the western border ($r = +0.37$), suggesting that KE advection from the west directly
modulates the lateral $K_e$ import into the LEC domain — consistent with the boundary
term being partly driven by the upper-level flux across the western domain boundary.

% ---------------------------------------------------------------
% Panel (d): PV 200
% ---------------------------------------------------------------

For the $\mathrm{PV}_{200}$ field
(Fig.~\ref{fig:9_pearson_epall_by_field_type}d), the overall strongest correlations in
the entire figure are found. Before discussing them, the correct physical interpretation
of the spatial features requires attention to the sign convention of potential vorticity
in the Southern Hemisphere: $\mathrm{PV}_{200}$ is \emph{always negative} here
($\approx -3$ to $-7$ PVU in the composite), since the Coriolis parameter is negative
and PV scales with $f$. More negative values therefore indicate \emph{more cyclonic}
circulation; less negative (i.e., higher) values correspond to anticyclonic conditions.
The EPALL composite shows a cyclonic anomaly to the \emph{west} of the surface cyclone
(anomaly $\approx -0.51$ PVU) and an anticyclonic anomaly to the \emph{east}
($\approx +0.40$ PVU), consistent with the well-known trough–ridge dipole embedded in
the Rossby wave train that characterizes South Atlantic extratropical cyclones.

The E–W contrast (east $-$ west) is therefore positive in the composite (east less
negative than west), and it reflects the degree of this trough-west / ridge-east
configuration. It correlates positively with $G_e$ ($r = +0.59$, $\rho = +0.57$),
$A_e$ ($r = +0.57$, $\rho = +0.59$), and $C_a$ ($r = +0.45$, $\rho = +0.47$): a
larger contrast — that is, a stronger cyclonic upper trough to the west combined with a
more pronounced anticyclonic ridge to the east — is associated with larger diabatic
generation, eddy APE, and baroclinic conversion. For $G_e$, the physical mechanism
likely involves the PV-induced ascending motion beneath the upper trough: by
invertibility \citep{hoskins1985use}, the negative PV anomaly to the west drives
asymmetric mid-tropospheric lifting, concentrated near the axis of maximum temperature
gradient between the trough and the downstream ridge. This lifting sustains
precipitation and latent heat release, thereby increasing $G_e$
\citep{attinger2021systematic}. The additional contribution of the downstream
anticyclonic ridge (less negative PV to the east) may further focus the ascending
branch by inhibiting lateral energy dispersion, so the correlation of $G_e$ with the
eastern-side features ($r_{\mathrm{border\_east}} = +0.46$, $r_{\mathrm{sector\_east}}
= +0.40$) is also consistent with this picture.

For $C_a$ and $A_e$, the strong negative correlations with the western border
($r_{C_a} = -0.30$, $r_{A_e} = -0.44$) and western sector ($r_{C_a} = -0.24$,
$r_{A_e} = -0.37$) indicate that a \emph{stronger} cyclonic PV anomaly to the west
(more negative $\mathrm{PV}_{200}$) is associated with \emph{larger} baroclinic
conversion and eddy APE. This is consistent with the PV-coupling mechanism for
baroclinic development \citep{hoskins1985use,davis1991potential}: the upper-level
cyclonic PV anomaly to the west of the surface low induces stretching and low-level
cyclonic circulation, reinforcing the surface pressure gradient and the temperature
contrast across the warm front. Cyclones in which this trough-to-west configuration is
more intense therefore exhibit higher $C_a$ and $A_e$. Likewise, the E-W contrast
correlation with $C_a$ ($r = +0.45$) captures the same mechanism from the
combined-contrast perspective.

We note that the correlation of the western border $\mathrm{PV}_{200}$ with
$G_e$ is not significant ($r = -0.05$, $\rho = -0.01$), indicating that the
cyclonic PV anomaly to the west, in isolation, is not a robust predictor of diabatic
APE generation at the population level. Instead, it is the full E–W contrast — the
coherent trough-west + ridge-east dipole — that best predicts $G_e$. This suggests
that the ascending branch favourable for latent heat release is determined less by the
absolute intensity of the western trough than by the phase relationship between the
trough and the downstream ridge.

The domain-mean $\mathrm{PV}_{200}$ does not appear above the $|r| = 0.20$ threshold
for the primary canonical terms $C_a$, $A_e$, $K_e$, or $C_k$, confirming that the
\emph{absolute} level of upper-level PV is less informative than the \emph{spatial
gradient} for predicting LEC term magnitudes. The marginal exception is $G_e$
($r = +0.25$ with the domain mean): a less negative (less cyclonic) overall domain PV
— driven by the anticyclonic eastern ridge dominating the spatial average — is
moderatly associated with more diabatic generation, consistent with the broader
trough–ridge-dipole narrative described above.

% ---------------------------------------------------------------
% Panel (e): PV 850
% ---------------------------------------------------------------

The $\mathrm{PV}_{850}$ panel is not discussed in detail because of its notably sparse
significant associations (7 of 91 pairs, 7.7\%; maximum $|r| = 0.27$). This sparsity
suggests that lower-tropospheric PV structure in the storm-centered frame is a
particularly noisy predictor of LEC term magnitudes at the population level, consistent
with the sensitivity of low-level PV to fine-scale, transient diabatic processes
(precipitation, surface fluxes) that vary widely from event to event. In contrast,
the upper-level PV (which has a strong and consistent E–W gradient across the EPALL
composite) provides the dominant PV signal in the correlation analysis.

% ---------------------------------------------------------------
% Synthesis
% ---------------------------------------------------------------

Taken together, the patterns in Fig.~\ref{fig:9_pearson_epall_by_field_type} reveal a
hierarchy of dynamical field informativeness for predicting LEC term magnitudes across
the full cyclone population: $\mathrm{AFC}_{250}$ and $\mathrm{KE\_adv}_{250}$ have
the highest fraction of pairs exceeding the 0.20 threshold (40\% and 38\%,
respectively), followed by $\mathrm{AdvT}_{850}$ (33\%) and $\mathrm{PV}_{200}$ (20\%).
However, the amplitudes of $\mathrm{AdvT}_{850}$ and $\mathrm{KE\_adv}_{250}$ are
themselves strongly correlated ($r_{\mathrm{abs}} = +0.61$), suggesting that these two
fields partly capture the same dynamical signal (the overall intensity of the
storm-relative advective forcing). In terms of individual correlation magnitude,
the E–W contrast of $\mathrm{PV}_{200}$ stands out as the leading single predictor
($|r|_{\max} = 0.59$ with $G_e$), followed by the northern-sector cold advection in
$\mathrm{AdvT}_{850}$ ($|r|_{\max} = 0.55$ with $A_e$) and the domain-amplitude AFC
($|r|_{\max} = 0.53$ with $K_e$).

Two overarching physical narratives emerge from this analysis. First, the negative
correlations of $C_a$ and $A_e$ with the \emph{western} and \emph{northern} sectors
of $\mathrm{AdvT}_{850}$ — combined with the positive correlations of $BA_e$ with
the \emph{eastern} features — reflect the canonical cold-front/warm-sector anatomy of
Southern Hemisphere extratropical cyclones (cold air to the west and northwest, warm
sector to the northeast), where a stronger frontal temperature contrast supports deeper
baroclinic development. Second, the dominant role of the E–W contrast in
$\mathrm{PV}_{200}$ — with cyclonic anomaly to the west and anticyclonic anomaly to
the east — points to the trough–ridge dipole of the Rossby wave train as the primary
upper-level energizing configuration, acting jointly through PV-induced ascent
(driving $G_e$) and baroclinic coupling between the upper trough and the surface low
(driving $C_a$ and $A_e$) \citep{hoskins1985use,davies2015quasigeostrophic,de2025lorenz}.
These findings are consistent with the composite structures shown in
Sect.~\ref{sec:composites}.

It should be noted that Fig.~\ref{fig:9_pearson_epall_by_field_type} shows
$|r|$ rather than signed $r$, which conceals the directionality of the associations.
Interpretation of the physical meaning of each entry therefore requires knowledge of
the sign, which is provided by the analysis underlying this section and is reported in
the supplementary material (Table~S?).

% ---------------------------------------------------------------
% Limitações — parágrafo opcional para inserção
% ---------------------------------------------------------------
%
% Possible limitations paragraph (insert as needed):
%
% Several limitations apply to the correlation analysis presented here.
% First, the use of Pearson $r$ assumes linear monotonic relationships;
% departures from linearity, while shown to be modest for most associations
% (Spearman $\rho$ agreed in sign and had similar magnitude for 100\% of
% displayed pairs), cannot be entirely ruled out. Second, the 13 spatial
% features used to summarize each 2D field are not orthogonal (e.g., the
% domain mean and sector means share information), so the number of
% effectively independent tests is lower than the nominal 455 tested pairs;
% however, the Bonferroni-corrected significance threshold at 455 tests
% ($|r| > 0.054$) remains well below the 0.20 display threshold.
% Third, while all shown correlations are formally highly significant
% ($p < 10^{-21}$), the effect sizes are modest ($|r| \leq 0.59$), and reflect
% population-level associations that do not necessarily predict individual
% cyclone behavior well. Fourth, correlation does not imply causality:
% the dynamical fields and LEC terms are computed from the same reanalysis
% data at overlapping time windows, so shared spatiotemporal covariance
% is expected; the physical interpretation provided here draws on additional
% theoretical and observational support but should be treated as hypothesis
% generation rather than mechanistic proof.

%% FIM DO TEXTO SUGERIDO %%
